% Source: physical PDF pp. 180--183 (printed pp. 157--160), from the Section
% 27.8 heading on physical p. 180 through Eq. (27.8.10) on physical p. 183,
% immediately before the Section 27.9 heading.
% Inventory: Eqs. (27.8.1)--(27.8.10), no unnumbered displays, one linked
% citation to Reference 11, two ordinary footnotes, and no figures, tables,
% or typographic dividers.
% Source anomaly preserved: the first lines of Eqs. (27.8.1) and (27.8.4)
% visibly contain a second \phi after the transformation spinor, unlike the
% cited Eq. (26.3.15).
% Uncertainties: none.

\subsection[Another Approach: Gauge-Invariant Supersymmetry Transformations]
{Another Approach: Gauge-Invariant Supersymmetry Transformations\footnotemark}
\footnotetext{This section lies somewhat out of the book's main line of
development and may be omitted in a first reading.}

It is somewhat disturbing that the supersymmetry transformation rules
discussed so far involve ordinary spacetime derivatives, not gauge-invariant
derivatives. For instance, in a \(U(1)\) gauge theory, the transformation of
the component fields of a chiral scalar superfield is given by
Eqs.~\eqref{eq:26.3.15}--\eqref{eq:26.3.17} as
\begin{equation}
  \begin{aligned}
    \delta\psi_\alpha
    &=-i\sqrt{2}\,\partial_\mu\phi\,
      (\sigma^\mu\bar\alpha)_\alpha\phi
      +\sqrt{2}\mathcal{F}\alpha_\alpha\,,\\
    \delta\mathcal{F}
    &=-i\sqrt{2}\,
      \bar\alpha\bar\sigma^\mu\partial_\mu\psi\,,\\
    \delta\phi
    &=\sqrt{2}\bigl(\alpha\psi\bigr)\,.
  \end{aligned}
  \tag{27.8.1}\label{eq:27.8.1}
\end{equation}
One might have thought that in the transformation of a chiral superfield
that carries \(U(1)\) charge \(q\), the ordinary spacetime derivatives in
Eq.~\eqref{eq:27.8.1} should be replaced with gauge-covariant derivatives,
given in terms of the \(U(1)\) gauge field \(V_\mu\) by
\begin{equation}
  D_\mu=\partial_\mu-iqV_\mu\,.
  \tag{27.8.2}\label{eq:27.8.2}
\end{equation}
With such gauge-invariant supersymmetry transformations of the chiral
superfields, one would still attempt to formulate supersymmetry
transformations of the gauge supermultiplet that involve only the physical
and auxiliary fields \(V_\mu\), \(\lambda\), and \(D\):
\begin{equation}
  \begin{aligned}
    \widetilde{\delta}V_\mu
    &=-i\left(
      \alpha\sigma_\mu\bar\lambda
      \bar\alpha\bar\sigma_\mu\lambda\right)\,,\\
    \widetilde{\delta}\lambda_\alpha
    &=f_{\mu\nu}(\sigma^{\mu\nu}\alpha)_\alpha
      +iD\alpha_\alpha\,,\\
    \widetilde{\delta}D
    &=\alpha\sigma^\mu\partial_\mu\bar\lambda
      -\bar\alpha\bar\sigma^\mu\partial_\mu\lambda\,,
  \end{aligned}
  \tag{27.8.3}\label{eq:27.8.3}
\end{equation}
in which ordinary spacetime derivatives appear because the gauge superfield
carries no \(U(1)\) charge.

This doesn't work. The algebra of these transformations does not close: the
commutator of two of the modified supersymmetry transformations is not a
linear combination of bosonic symmetry transformations, such as spacetime
translations and gauge transformations. It follows that it is not possible
to construct a Lagrangian for the chiral and gauge superfields that would be
invariant under these modified supersymmetry transformations, because if
there were such a Lagrangian then it would have to be also invariant under
the commutators of these transformations, so that these commutators would
have to be bosonic symmetries of the Lagrangian.

In 1973 de Wit and Freedman~\hyperref[ch27-ref-11]{[11]} showed that the
supersymmetry algebra could be made to close by modifying the supersymmetry
transformation properties of the chiral superfields not only by changing
ordinary derivatives to gauge-invariant derivatives, but by also adding an
extra term in the transformation of the \(\mathcal{F}\)-component, so that
for \(U(1)\) gauge theories the modified supersymmetry transformation rules
read
\begin{equation}
  \begin{aligned}
    \widetilde{\delta}\psi_\alpha
    &=-i\sqrt{2}\,D_\mu\phi\,
      (\sigma^\mu\bar\alpha)_\alpha\phi
      +\sqrt{2}\mathcal{F}\alpha_\alpha\,,\\
    \widetilde{\delta}\mathcal{F}
    &=-i\sqrt{2}\bar\alpha\bar\sigma^\mu D_\mu\psi
      -2iq\phi\bigl(\bar\alpha\bar\lambda\bigr)\,,\\
    \widetilde{\delta}\phi
    &=\sqrt{2}\bigl(\alpha\psi\bigr)\,.
  \end{aligned}
  \tag{27.8.4}\label{eq:27.8.4}
\end{equation}
With this change, they were also able to construct a Lagrangian that is
invariant under the transformations \eqref{eq:27.8.3}--\eqref{eq:27.8.4}
and that turned out to be just the one we have found in Sections 27.1 and
27.2.

There is nothing wrong with continuing to use the conventional
transformation rules \eqref{eq:27.8.1}, so we do not need the de
Wit--Freedman formalism to deal with supersymmetric gauge theories.
Nevertheless this formalism is of some interest, because in supergravity
theories the analog of the conventional formalism is very cumbersome. As
described in Chapter 31, the formalism that has been chiefly used to derive
physically interesting results in supergravity theories follows an approach
like that of de Wit and Freedman, with supersymmetry transformation rules
involving covariant derivatives in place of ordinary derivatives, rather
than an approach based on conventional supersymmetry transformations like
those of Eq.~\eqref{eq:27.8.1}. It is therefore of some interest to
understand the relation between the de Wit--Freedman formalism and the
conventional approach in the relatively simple context of \(U(1)\) gauge
theory, and in particular to explain the origin of the extra term in the
transformation rule for \(\mathcal{F}\).

In writing supersymmetry transformations \eqref{eq:27.8.3} that did not
involve the components \(C\), \(M\), \(N\), or \(\omega\) of the gauge
superfield \(V\), de Wit and Freedman were implicitly adopting the
Wess--Zumino gauge discussed in Section 27.1. But the choice of Wess--Zumino
gauge is not invariant under either the conventional supersymmetry
transformations \eqref{eq:26.2.11}--\eqref{eq:26.2.17} or the extended gauge
transformations \eqref{eq:27.1.17}, so once we adopt this gauge both
symmetries are lost. We can, however, define a \emph{combined}
transformation, acting on fields in Wess--Zumino gauge, which consists of a
conventional supersymmetry transformation followed by an extended gauge
transformation that takes us back to Wess--Zumino gauge. \emph{This is the
de Wit--Freedman transformation \(\widetilde{\delta}\)}\footnote{This was not
shown explicitly by de Wit and Freedman. But, in fact, although the point of
their paper was to emphasize that the details of the transformations
\eqref{eq:27.8.3} and \eqref{eq:27.8.4} could be inferred from the
requirement of a closed supersymmetry algebra (also for non-Abelian gauge
theories), they remarked that they had actually found these transformations
by identifying the fermionic transformations that survive in Wess--Zumino
gauge.}.

To construct the de Wit--Freedman transformations in this way, note that for
a gauge superfield that satisfies the Wess--Zumino gauge conditions
\(C=M=N=\omega=0\), the transformation rules
\eqref{eq:26.2.11}--\eqref{eq:26.2.14} give
\begin{equation}
  \delta C=0,\qquad
  \delta\omega_\alpha=-i(V_\mu\sigma^\mu\bar\alpha)_\alpha,\qquad
  \delta M=-\bigl(\alpha\lambda+\bar\alpha\bar\lambda\bigr),\qquad
  \delta N=i\bigl(\alpha\lambda-\bar\alpha\bar\lambda\bigr)\,.
  \tag{27.8.5}\label{eq:27.8.5}
\end{equation}
According to Eqs.~\eqref{eq:27.1.17}, we can get back to Wess--Zumino gauge
by performing an infinitesimal extended gauge transformation
\eqref{eq:27.1.13}:
\begin{equation}
  V\longrightarrow V+\frac{i}{2}\bigl[\Omega-\Omega^*\bigr]\,,
  \tag{27.8.6}\label{eq:27.8.6}
\end{equation}
where \(\Omega\) is a left-chiral superfield with components
\begin{equation}
  \phi^\Omega=0,\qquad
  \psi_\alpha^\Omega
    =i\sqrt{2}(V_\mu\sigma^\mu\bar\alpha)_\alpha,\qquad
  \mathcal{F}^\Omega
  =-2\bigl(\bar\alpha\bar\lambda\bigr)\,.
  \tag{27.8.7}\label{eq:27.8.7}
\end{equation}
According to Eq.~\eqref{eq:27.1.11}, this extended gauge transformation
induces on a chiral superfield of charge \(q\) the transformation
\begin{equation}
  \delta'\Phi=iq\Omega\Phi\,.
  \tag{27.8.8}\label{eq:27.8.8}
\end{equation}
Using the multiplication rules
\eqref{eq:26.3.27}--\eqref{eq:26.3.29}, the transformation of the components
of \(\Phi\) is
\begin{equation}
  \begin{aligned}
    \delta'\psi_\alpha
    &=-\sqrt{2}q\phi\,
      (V_\mu\sigma^\mu\bar\alpha)_\alpha\,,\\
    \delta'\mathcal{F}
    &=-2iq\phi\bigl(\bar\alpha\bar\lambda\bigr)
      -\sqrt{2}q\bar\alpha\bar\sigma^\mu V_\mu\psi\,,\\
    \delta'\phi&=0\,.
  \end{aligned}
  \tag{27.8.9}\label{eq:27.8.9}
\end{equation}
Adding this to Eq.~\eqref{eq:27.8.1} and comparing with
Eq.~\eqref{eq:27.8.4} shows that the de Wit--Freedman transformation is
indeed the combination of a conventional supersymmetry transformation and
the corresponding extended gauge transformation \eqref{eq:27.8.8}:
\begin{equation}
  \widetilde{\delta}\Phi=\delta\Phi+\delta'\Phi\,.
  \tag{27.8.10}\label{eq:27.8.10}
\end{equation}
